(1) 模型方程为 \(m\dot{\boldsymbol v}+m\gamma\boldsymbol v=-e\boldsymbol E(t)\)。取复振幅 \(v(t)=\Re[v_0e^{-i\omega t}]\)，有 \(m(\gamma-i\omega)v_0=-eE_0\)，所以
\[
v_0=-\frac{eE_0}{m(\gamma-i\omega)},\qquad
\boldsymbol J_0=-ne\boldsymbol v_0
=\sigma(\omega)\boldsymbol E_0,\quad
\sigma(\omega)=\frac{ne^2}{m(\gamma-i\omega)}.
\]

(2) 将导电电流纳入总位移电流：\(\nabla\times\boldsymbol H=\sigma E-i\omega\epsilon_0E=-i\omega\epsilon_0\epsilon_r E\)。故 \(\epsilon_r=1+i\sigma/(\epsilon_0\omega)\)，定义 \(\omega_p^2=ne^2/(m\epsilon_0)\)，化简得
\[
\epsilon_r(\omega)=1-\frac{\omega_p^2}{\omega(\omega+i\gamma)}.
\]
这里采用 \(e^{-i\omega t}\) 时间惯例，因此被动介质应有 \(\Im\epsilon_r>0\)。

(3) 复幅值的平均功率密度为 \(\overline p=\frac12\Re(\boldsymbol J_0\cdot\boldsymbol E_0^*)=\frac12\Re\sigma\,|E_0|^2\)。而
\[
\Re\sigma=\frac{ne^2\gamma}{m(\gamma^2+\omega^2)}>0,\qquad
\overline p=\frac{ne^2\gamma|E_0|^2}{2m(\gamma^2+\omega^2)}.
\]
对运动方程乘速度再作周期平均也得 \(\overline{ne\boldsymbol E\cdot(-\boldsymbol v)}=nm\gamma\overline{v^2}\)，从微观阻尼角度核验正耗散。

(4) 在 \(\omega\ll\gamma\) 时 \(\sigma\simeq ne^2/(m\gamma)\)，恢复欧姆电导；在 \(\omega\gg\gamma\) 时 \(\sigma\simeq i ne^2/(m\omega)\)，主要为无耗散的感性响应。固定 \(\omega \) 后令 \(\gamma \to 0^{+}\)，\(\epsilon_r\to1-\omega_p^2/\omega^2\)，当 \(0<\omega<\omega_p\) 时介电常数实部为负，均匀无碰撞等离子体中的电磁平面波不能以实波数传播。该结论限于忽略热速度、边界与碰撞的简单 Drude 模型。
